\documentclass[10pt,twocolumn]{article}
\usepackage[T1]{fontenc}
\usepackage[utf8]{inputenc}
\usepackage{lmodern}
\usepackage[margin=1.8cm,top=2.4cm,bottom=2cm,columnsep=0.6cm]{geometry}
\usepackage{fancyhdr}
\usepackage{titlesec}
\usepackage{booktabs}
\usepackage{amsmath,amssymb,amsfonts}
\usepackage{microtype}
\usepackage{xcolor}
\usepackage{mdframed}
\usepackage{parskip}
\usepackage{enumitem}
\usepackage{hyperref}

\definecolor{rulered}{RGB}{180,30,30}
\definecolor{boxgray}{RGB}{245,245,245}
\definecolor{keywordgray}{RGB}{100,100,100}

\pagestyle{fancy}
\fancyhf{}
\fancyhead[L]{\textsc{\small Claude Mag}}
\fancyhead[C]{\small\textit{Machine Learning Research Digest}}
\fancyhead[R]{\small 2026-11-23}
\fancyfoot[C]{\small\thepage}
\renewcommand{\headrulewidth}{0.8pt}

\titleformat{\section}{\normalsize\bfseries\color{rulered}}{}{0pt}{}%
  [\vspace{-4pt}\color{rulered}\rule{\columnwidth}{0.4pt}]
\titlespacing{\section}{0pt}{8pt}{4pt}

\hypersetup{colorlinks=true,urlcolor=rulered,linkcolor=black}

\begin{document}
\setlength{\parindent}{0pt}
\setlength{\parskip}{4pt}

{\small\color{keywordgray}\textbf{Keywords:}
  Multimodal Reasoning \textbullet{} Visual Primitives \textbullet{} Spatial Grounding \textbullet{} DeepSeek \textbullet{} Reference Gap}\par\vspace{4pt}

{\LARGE\bfseries\raggedright Thinking with Visual Primitives}\par\vspace{4pt}
{\large\itshape\raggedright DeepSeek treats spatial markers as minimal units of thought, closing the Reference Gap in multimodal reasoning at 98.7\% spatial accuracy}\par\vspace{6pt}
{\small\textbf{By} Ruijie Lu, Yiyang Ma, Gang Zeng, Ting Chen, et al.\ (DeepSeek-AI; Peking University; Tsinghua University)}\par\vspace{2pt}
{\small\color{keywordgray}NeurIPS 2026 Poster \textbullet{} Technical Report 30 Apr 2026}
\par\vspace{2pt}\noindent{\color{rulered}\rule{\columnwidth}{0.6pt}}\vspace{6pt}

Lu, Ma, Zeng, Chen, and colleagues from DeepSeek-AI, Peking University, and Tsinghua University propose treating spatial markers---points and bounding boxes---as minimal units of thought that a multimodal model generates auto-regressively alongside text tokens.
Interleaving these \emph{visual primitives} into the chain of thought resolves what the authors call the Reference Gap: the ambiguity between language tokens and precise spatial locations in cluttered visual scenes.
The resulting model achieves 98.7\% accuracy on a spatial reasoning benchmark, 1--2 percentage points above the strongest competing systems.

\begin{mdframed}[backgroundcolor=boxgray,linecolor=rulered,linewidth=1pt,
    innerleftmargin=6pt,innerrightmargin=6pt,
    innertopmargin=6pt,innerbottommargin=6pt]
\textbf{\small AT A GLANCE}\par\vspace{4pt}
\begin{description}[leftmargin=0pt,labelsep=4pt,itemsep=2pt,parsep=0pt,
    font=\small\bfseries]
  \item[Problem:] Language tokens cannot precisely refer to spatial locations in high-resolution or cluttered images; image cropping introduces cascading errors.
  \item[Why it matters:] Multi-step spatial reasoning is load-bearing for real-world visual AI; closing the Reference Gap enables a qualitative capability jump.
  \item[Method:] Point and box tokens in the vocabulary; visual encoder called on each primitive token at generation time; RLVR fine-tuning with spatial verifiers.
  \item[Result:] 98.7\% on spatial reasoning benchmark; 1--2 pp above GPT-5-V and Gemini 2.5 Ultra.
  \item[Evidence:] DeepSeek-V4-Flash MoE backbone; KV compression 4:1; 2,400-question in-house benchmark.
\end{description}
\end{mdframed}

\section{The Big Picture}

Multimodal language models have made rapid progress on image question-answering benchmarks, but dense spatial reasoning tasks --- tasks requiring multi-step localization in cluttered, high-resolution images --- remain a systematic weakness.
Existing methods work around this by cropping patches and feeding them as additional images.
This paper argues that cropping is the wrong abstraction level and that the solution belongs inside the model's reasoning chain.

The underlying diagnosis is what the authors call the \emph{Reference Gap}: language tokens denote concepts at a semantic level, but spatial queries require sub-semantic precision.
When a chain of thought says ``the second blue sphere in the top row,'' this phrase must ultimately resolve to a specific pixel region --- but the resolution step happens implicitly and unreliably.
Visual primitives make the resolution step explicit and verifiable.

\section{Why Existing Approaches Fall Short}

\textbf{Image cropping} requires the model to decide where to crop before reasoning about the crop, creating a circular dependency: the spatial reasoning required to crop the right region is the same capability the crop is supposed to provide.
Errors in the crop decision cascade into all subsequent reasoning steps.

\textbf{High-resolution encoders} increase the token budget for visual features but do not change the representational mismatch between language tokens and spatial coordinates.
More visual tokens do not close the Reference Gap; they only make it slightly narrower.

\textbf{Grounding-then-answering pipelines} run a visual grounding module as a preprocessing step and feed its outputs to the reasoning model.
Module errors are not correctable inside the reasoning chain; the system cannot iteratively refine its spatial hypotheses.

Thinking with Visual Primitives closes the gap by making spatial localization a first-class operation in the model's generation vocabulary, subject to in-context correction at any step.

\section{Core Mechanism}

The model's vocabulary is extended with two categories of visual primitive tokens:

\begin{itemize}[leftmargin=1em,itemsep=2pt]
  \item \textbf{Point tokens:} \texttt{<point x="0.42" y="0.67"/>} --- normalized image coordinates.
  \item \textbf{Box tokens:} \texttt{<box x1="0.21" y1="0.33" x2="0.58" y2="0.71"/>} --- bounding box specifying a spatial region.
\end{itemize}

Both are generated auto-regressively alongside text tokens.
When the model generates a primitive token, the visual encoder extracts features from the corresponding location or region and injects them into the next-token prediction context.
The model can generate, reason about, and revise spatial references entirely within the generation loop.

\section{Technical Formulation}

\subsection*{Mixed-Token Generation}

The model generates a sequence of text and primitive tokens:
\begin{equation}
y = (y_1, \ldots, y_T),\quad y_i \in \mathcal{V}_{\mathrm{text}} \cup \mathcal{V}_{\mathrm{prim}}
\label{eq:gen}
\end{equation}

At step $i$, if $y_i \in \mathcal{V}_{\mathrm{prim}}$ (a spatial token), the visual encoder $\phi_v$ extracts a feature vector from the corresponding region:
\begin{equation}
v_i = \phi_v\!\left(I,\; \mathrm{region}(y_i)\right)
\label{eq:feat}
\end{equation}

This feature vector is concatenated with the standard context embedding $c_i$ before computing the next-token distribution:
\begin{equation}
P(y_{i+1} \mid y_{\leq i}) = \mathrm{softmax}\!\left(W [c_i;\, v_i]\right)
\label{eq:ntp}
\end{equation}

For text tokens, $v_i = \mathbf{0}$ and standard language modeling applies.

\subsection*{KV-Cache Compression}

The DeepSeek-V4-Flash backbone compresses every 4 consecutive visual token key-value pairs into a single cached entry:
\begin{equation}
\mathrm{KV}_{\mathrm{vis},k} = \mathrm{Compress}\!\left(\mathrm{KV}_{4k}, \mathrm{KV}_{4k+1}, \mathrm{KV}_{4k+2}, \mathrm{KV}_{4k+3}\right)
\label{eq:kv}
\end{equation}

This reduces visual memory cost by $4\times$, enabling dense visual reasoning over long chains of thought without exceeding memory limits.

\section{Learning or Inference Procedure}

Training has three phases:

\begin{enumerate}[leftmargin=1.5em,itemsep=1pt]
  \item \textbf{Supervised pretraining on spatial annotation:} Train the model to generate correct primitive tokens given ground-truth spatial annotations (e.g., object detection data).
  \item \textbf{Chain-of-thought reasoning fine-tuning:} Train on spatial reasoning problems with multi-step chains of thought that interleave primitive tokens.
  \item \textbf{RLVR fine-tuning:} Refine with verifiable reward signals---correctness of spatial reasoning outputs checked by a deterministic spatial verifier---using a GRPO-style algorithm.
\end{enumerate}

At inference, the model generates normally.
The primitive-based reasoning mode activates when the prompt contains explicit spatial reasoning trigger keywords.

\section{What the Guarantee Says}

There is no formal guarantee.
The empirical claim is that the locate--mark--read loop enabled by visual primitives provides a systematic advantage on tasks requiring multi-step spatial reference, over both language-only and crop-based approaches.

A known limitation: the primitive mode currently requires explicit trigger keywords.
Extending automatic triggering --- deciding when to use primitives without explicit prompting --- is listed as future work.

\section{Experimental Findings}

The spatial reasoning benchmark contains 2,400 questions requiring multi-step spatial reasoning over cluttered scenes.

\begin{center}
\small
\begin{tabular}{lr}
\toprule
\textbf{Model} & \textbf{Accuracy} \\
\midrule
GPT-5-V (image cropping) & 97.1\% \\
Gemini 2.5 Ultra & 96.8\% \\
Claude Muse-V & 96.4\% \\
Text-only reasoning (no primitives) & 89.3\% \\
Image cropping (ablation) & 93.7\% \\
Thinking with Visual Primitives & \textbf{98.7\%} \\
\bottomrule
\end{tabular}
\end{center}

Performance on standard multimodal benchmarks (MMBench, SeedBench, MMStar) is maintained or marginally improved relative to the base DeepSeek-V4-Flash model.

\section{Ablations and Interpretation}

\begin{itemize}[leftmargin=1em,itemsep=2pt]
  \item \textbf{Without RLVR fine-tuning:} Accuracy drops from 98.7\% to 94.1\%, confirming that reward-based refinement is essential for precision.
  \item \textbf{Without chain-of-thought reasoning fine-tuning:} 91.8\%, showing that multi-step reasoning training contributes independently of the RLVR phase.
  \item \textbf{Point tokens only (no box tokens):} 96.3\%, a 2.4~pp drop, confirming that region-level primitives add value over point localization alone.
  \item \textbf{Trigger keywords:} Accuracy drops to 72.1\% when primitive tokens are always generated (no trigger), because the model generates spurious primitives that distract from text-based reasoning steps.
    Targeted triggering is a critical design choice.
\end{itemize}

\section*{Reference}
R.\ Lu, Y.\ Ma, G.\ Zeng, T.\ Chen, et al.
\textbf{Thinking with Visual Primitives.}
DeepSeek-AI Technical Report. NeurIPS 2026 Poster.
\url{https://neurips.cc/virtual/2026/poster/154630}

\end{document}
